\section{Experiments}
\label{sec:experiments}

Our studies are organized by claim. Final CIFAR-10 claims require
50k generated samples for FID, NFE coverage that includes 10 and 20, and more
than one seed for reliability language, with matched base, native-coordinate
linear, Karras/EDM-style \citep{karras2022edm}, project-owned offline-proxy,
and our schedules. The retained CIFAR backend is PNDM
\citep{liu2022pndm} with the Euler solver. Text-to-image diagnostics require
matched prompt, seed, model, solver, NFE, and guidance-scale keys across base,
published AYS \citep{sabour2024ays}, and our schedules. The repository
also contains STORK adapter tests and scheduler code \citep{tan2025stork}, but
this draft does not report paper-grade STORK quality results.

\begin{table}[t]
  \centering
  \caption{Real PNDM/CIFAR-10 smoke gate for our implementation. The run uses
  Euler, NFE 4, seed 0, and 32 generated samples per
  schedule. FID is reported only to verify that scoring produced finite,
  non-empty outputs; it is not paper-grade evidence.}
  \label{tab:smoke-gate}
  \begin{tabular}{lrrrr}
    \toprule
    Schedule & NFE & Seed & Samples & FID \\
    \midrule
    Base & 4 & 0 & 32 & 295.77 \\
    ours & 4 & 0 & 32 & 258.82 \\
    \bottomrule
  \end{tabular}
\end{table}


Table~\ref{tab:smoke-gate} is a pipeline gate, not a paper result. It verifies
that the repository can materialize our schedule, generate samples for
base and our schedule, compute non-empty FID rows, and write calibration-cost
metadata. The 32-sample FID values cannot support any image-quality claim.

\begin{table}[t]
  \centering
  \caption{PNDM/CIFAR-10 50k FID for the Euler solver. Rows report
  mean $\pm$ standard error over three seeds (0, 1, 2), base, linear
  native-coordinate, fixed Karras/EDM-style, authorized project-owned
  offline, and our schedules. FID is lower-is-better and each
  seed uses 50k generated images. The final two columns are paired
  ours-minus-baseline FID gaps; negative values favor ours.}
  \label{tab:cifar10-pndm-euler-50k-fid-seeds0-1-2}
  {\scriptsize
  \setlength{\tabcolsep}{2pt}
  \resizebox{\linewidth}{!}{%
  \begin{tabular}{rrrrrrrr}
    \toprule
    NFE & base & linear & Karras & offline & ours &
    ours $-$ offline & ours $-$ Karras \\
    \midrule
    10 & 20.57 $\pm$ 0.08 & 20.57 $\pm$ 0.08 & 25.84 $\pm$ 0.10 & 19.20 $\pm$ 0.09 & 17.54 $\pm$ 0.08 & -1.659 $\pm$ 0.017 & -8.302 $\pm$ 0.031 \\
    20 & 11.17 $\pm$ 0.07 & 11.20 $\pm$ 0.07 & 10.64 $\pm$ 0.09 & 10.88 $\pm$ 0.06 & 10.71 $\pm$ 0.07 & -0.175 $\pm$ 0.059 & 0.072 $\pm$ 0.024 \\
    \bottomrule
  \end{tabular}
  }
  }
\end{table}


Table~\ref{tab:cifar10-pndm-euler-50k-fid-seeds0-1-2} records the current
three-seed 50k CIFAR-10 run for NFE 10 and 20, including native-coordinate
linear, fixed Karras/EDM-style, and authorized project-owned offline-proxy
baselines. The offline-proxy rows come from the project-owned coordinate-search optimizer; they
are not published AYS schedules. The result supports only a bounded
controlled-benchmark statement
for the PNDM/CIFAR-10 backend, Euler solver, these two NFE values, and the
reported baseline set. Appendix
Figure~\ref{fig:cifar10-pndm-euler-50k-seeds0-1-2} shows the same aggregate
with per-seed points. Ours improves FID over base, linear, and offline at
both NFE values. The Karras-style grid is worse than ours at NFE 10 but
slightly better at NFE 20, so the positive claim is scoped away from uniform
superiority over fixed schedules. Appendix
Figure~\ref{fig:cifar10-pndm-euler-schedule-profile-nfe20} visualizes the
NFE-20 schedule profiles for seed 0. Appendix
Table~\ref{tab:cifar10-pndm-euler-50k-reuse-seed0-schedule} gives a narrow
actual reuse check for the same backend: reusing the seed-0 schedule
for seeds 1 and 2 has zero observed FID gap because the continuous schedules
round to the same PNDM Euler integer timesteps. Appendix
Tables~\ref{tab:sd15-euler-nfe10-cfg-sweep-pairwise} and
\ref{tab:sdxl-euler-nfe10-cfg-sweep-pairwise} record the matched
text-to-image CFG sweeps where published AYS schedules enter the retained
experiments. SD1.5 is a negative case for our schedule under CLIPScore and
ImageReward. SDXL is condition-dependent: at CFG 5.0 our schedule improves
ImageReward over base but trails AYS, while at CFG 7.5 it improves the mean
CLIPScore and ImageReward deltas against both base and AYS. These sweeps
support CFG-matched coverage and boundary analysis; broad text-to-image
quality improvement remains unsupported. Appendix
Figure~\ref{fig:sd15-euler-nfe10-cfg7p5-schedule-profile} shows the
representative base/AYS/ours sigma-grid diagnostic for CFG 7.5 and seed 0.

\begin{table}[t]
  \centering
  \caption{Calibration cost accounting for current paper-facing runs.
  The quality column reports base-minus-ours FID for CIFAR-10 and
  ours-minus-base ImageReward for text-to-image, so positive values favor
  ours. Model-evaluation equivalents count retained calibration
  cost and the evaluated generation batch, excluding unrelated pilot
  experiments.}
  \label{tab:calibration-cost-summary}
  \begin{tabular}{lllrrr}
    \toprule
    Setting & condition & quality vs. base & cal. evals (M) & gen. evals (M) & cal./gen. \\
    \midrule
    CIFAR-10 & NFE 10 & +3.026 $\pm$ 0.019 FID & 12.73 & 1.500 & 8.49$\times$ \\
    CIFAR-10 & NFE 20 & +0.464 $\pm$ 0.005 FID & 12.73 & 3.000 & 4.24$\times$ \\
    SD1.5 & CFG 5 & -0.133 IR & 25.46 & 0.003 & 8{,}487$\times$ \\
    SD1.5 & CFG 7.5 & -0.084 IR & 25.46 & 0.003 & 8{,}487$\times$ \\
    SD1.5 & CFG 10 & -0.087 IR & 25.46 & 0.003 & 8{,}487$\times$ \\
    SDXL & CFG 5 & +0.079 IR & 0.21 & 0.003 & 68.62$\times$ \\
    SDXL & CFG 7.5 & +0.122 IR & 0.21 & 0.003 & 68.62$\times$ \\
    \bottomrule
  \end{tabular}
\end{table}


Table~\ref{tab:calibration-cost-summary} reports the current calibration-cost
accounting. The CIFAR-10 gain over base comes with calibration cost that exceeds
the evaluated generation batch by 4.24--8.49 times. The SD1.5 sweep has a much
larger calibration-to-generation ratio for the small matched prompt batch and
negative ImageReward deltas. The SDXL rows use a deliberately lighter
calibration configuration and have lower cost ratios, but they still require
schedule reuse across many matched prompt batches for amortization. These
numbers support cost-aware reporting and leave deployment efficiency
unsupported.

\textbf{Takeaway.} The three-seed CIFAR setting improves over base,
native-linear, and authorized project-owned offline-proxy schedules at NFE 10
and 20. The Karras comparison is mixed and stronger at NFE 20. The matched
SD1.5 and SDXL sweeps establish text-to-image coverage, but their automated
metrics are condition-dependent and constrain the claims.
